\section{Synthesis-Aided Query Optimization}
\label{slab:sec:alg}


\revisioncancut{Section~\ref{slab:sec:toplevelalg} first presents a high-level overview of our synthesis-aided query optimization algorithm. 
Subsequent sections present the technical details of each component. 
We define dataflows for SQL queries --- which is one of our key contributions --- and use that to develop a monotonic scoring in Section~\ref{slab:sec:scorefunction}.
Section~\ref{slab:sec:prioritizingsearch} presents our query search algorithm. 
Finally, we describe how we check for equivalence and rank queries in  Sections~\ref{slab:sec:checkingequiv} and~\ref{slab:sec:perfranking}, respectively.}

 
\newcounter{mycounter}


\subsection{Top-Level Algorithm}
\label{slab:sec:toplevelalg}
Our top-level algorithm is shown in Algorithm~\ref{slab:alg:toplevelalg}. It accepts as input a query $\inputquery$ and an integrity constraint $\integrityconstraint$, and returns an output query $\query$ that is semantically equivalent to and faster than $\inputquery$. 
%In other words, if $\query$ passes our equivalence checker (line~\ref{slab:alg:toplevelalg:checkequiv}) and is ranked as faster than $\inputquery$ (line~\ref{slab:alg:toplevelalg:checkfaster}), then we return $\query$ at line~\ref{slab:alg:toplevelalg:returnquery}\footnote{For simplicity of our notation, here we only return the fastest query but one can also return the entire ranked list of rewritten queries to the user.}.
$\query$ is $\bot$ if it does not find one such query within the given time limit. 




Let us explain Algorithm~\ref{slab:alg:toplevelalg} in more detail. At a high-level, our algorithm is a \emph{new instantiation} of the counterexample-guided inductive synthesis (CEGIS) paradigm~\cite{solar2008program} --- which has found great success in the programming languages community --- for our query optimization problem. 
Our approach consists of (a) an inductive synthesizer (i.e., our query synthesizer) that aims to find a query $\query$ which satisfies a set $\counterexamples$ of \emph{counterexamples} and (b) a query checker (i.e., our equivalence checker) which checks whether or not $\query$ is equivalent to $\inputquery$ and generates a counterexample $\counterexample$ if not. 
One key advantage of using CEGIS is to enable efficient checking: checking a candidate query against $\counterexamples$ (line~\ref{slab:alg:toplevelalg:satisfycounterexample}) is typically significantly cheaper than calling an equivalence checker (line~\ref{slab:alg:toplevelalg:checkequiv}, e.g., \spes~\cite{zhou2022spes}). In other words, verification is used parsimoniously only when necessary. 
In our context, a counterexample $\counterexample$ is an input DB together with the desired output table returned by $\inputquery$. 




{In addition to using counterexamples, \toolname also uses a novel $\scorefunction$ function that defines dataflows for queries in order to guide the query synthesis process. 
While we defer the formal definition of $\scorefunction$ to Section~\ref{slab:sec:scorefunction}, from a high-level, it assigns a non-positive integer score\footnote{While \emph{non-positive} scores may seem counter-intuitive, they can be viewed as the negation of the cost, i.e., the higher the score, the lower the cost.} to each query $\query$ such that,  
$\query$ with a higher score is more likely to be an optimized query (i.e., semantically equivalent to and faster than $\inputquery$) and vice versa. 
Initially, $\counterexamples$ has no counterexamples and $\queryset$ is a empty set that stores equivalent queries (line~\ref{slab:alg:toplevelalg:init}). 
Then, we enter a CEGIS loop (lines~\ref{slab:alg:toplevelalg:whilebegin}--\ref{slab:alg:toplevelalg:whileend}). At line~\ref{slab:alg:toplevelalg:forloop}, it \emph{lazily} enumerates all queries in non-ascending order of their scores. In other words, $\toolname$ prioritizes the search for candidate queries $\query$ that are likely to optimize $\inputquery$ (i.e., with a higher score). 
Line~\ref{slab:alg:toplevelalg:satisfycounterexample} checks if $\query$ meets $\counterexamples$: if not, it continues to the next candidate. 
If $\query$ passes $\counterexamples$, we invoke the equivalence checker (line~\ref{slab:alg:toplevelalg:checkequiv}) to see if $\query$ is equivalent --- it returns a new counterexample $\counterexample$ if not equivalent, or returns $\emph{null}$ otherwise. 
For the latter case, line~\ref{slab:alg:toplevelalg:nocounterexample} adds $\query$ to $\queryset$; for the former, line~\ref{slab:alg:toplevelalg:addcounterexample} adds  $\counterexample$ to $\counterexamples$.
Upon termination, we rank queries in $\queryset$ (line~\ref{slab:alg:toplevelalg:checkfaster}) and return a ``best'' query that we believe can \emph{optimize} $\inputquery$.}



\begin{figure}[!t]
\vspace{-5pt}
\footnotesize
\centering

\begin{algorithm}[H]
\begin{algorithmic}[1]

\myprocedure{Optimize}{$\inputquery, \integrityconstraint$}

\Statex\Input{An input query $\inputquery$ and an integrity constraint $\integrityconstraint$.}


\Statex\Output{An equivalent and faster output query $\query$.}

\vspace{2pt}

\State $\counterexamples \assign \emptyset$; \  $\queryset \assign \emptyset$; 
\label{slab:alg:toplevelalg:init}



\While{not timed out}
\label{slab:alg:toplevelalg:whilebegin} 



\ForAll{$\query \in \textsc{LazyEnumerate}(\inputquery)$} \label{slab:alg:toplevelalg:forloop}
\Comment{{{\scriptsize \Circled{1}}} Query synthesizer.}


\State \textbf{if} $\query$ doesn't pass  $\counterexamples$ \textbf{then} \textbf{continue}; \label{slab:alg:toplevelalg:satisfycounterexample}
\Comment{Check against counterexs.}


\State $\counterexample \assign \textsc{CheckEquivalence}(\query, \inputquery, \integrityconstraint)$;\label{slab:alg:toplevelalg:checkequiv}
\Comment{{{\scriptsize \Circled{2}}} Equivalence checker.}

\vspace{1pt}

\If{$\counterexample = \emph{null}$} 
$\queryset.\emph{add}(\query)$
\label{slab:alg:toplevelalg:nocounterexample}
\Comment{$\emph{null}$ means $\query$ is equivalent to $\inputquery$.}

\Else \ $\counterexamples.\emph{add}( \counterexample )$;
\Comment{Otherwise, $\counterexample$ is a new counterexample.}
\label{slab:alg:toplevelalg:addcounterexample}

\EndIf

\EndFor


\EndWhile
\label{slab:alg:toplevelalg:whileend}

\vspace{-1pt}
\State $\query \assign \textsc{RankPerformance}(\queryset, \inputquery)$;
\label{slab:alg:toplevelalg:checkfaster}
\Comment{{{\scriptsize \Circled{3}}} Performance Ranker.}

\State \Return $\query$;
\label{slab:alg:toplevelalg:returnquery}

\caption{{Top-level algorithm of \toolname.}}
\label{slab:alg:toplevelalg}
\end{algorithmic}
\end{algorithm}
\vspace{-10pt}
\end{figure}





\begin{comment}
    
\begin{figure}[!t]
\small 
\centering

\begin{algorithm}[H]
\begin{algorithmic}[1]

\myprocedure{Optimize}{$\inputquery, \integrityconstraint$}

\Statex\Input{An input query $\inputquery$ and an integrity constraint $\integrityconstraint$.}


\Statex\Output{An equivalent and faster output query $\query$.}

\vspace{2pt}

\State $\counterexamples \assign \emptyset$; \  $\cost = 0; \ \queryset \assign \emptyset$; 
\label{slab:alg:toplevelalg:init}



\While{not timed out}
\label{slab:alg:toplevelalg:whilenottimeout} 



\ForAll{$\query \in \textsc{LazyEnumerate}(\cost)$} \label{slab:alg:toplevelalg:forloop}
\Comment{{{\scriptsize \Circled{1}}} Query synthesizer.}


\State \textbf{if} $\query$ doesn't pass  $\counterexamples$ \textbf{then} \textbf{continue}; \label{slab:alg:toplevelalg:satisfycounterexample}
\Comment{Check against counterexs.}


\State $\counterexample \assign \textsc{CheckEquivalence}(\query, \inputquery, \integrityconstraint)$;\label{slab:alg:toplevelalg:checkequiv}
\Comment{{{\scriptsize \Circled{2}}} Equivalence checker.}

\vspace{1pt}

\If{$\counterexample = \emph{null}$} 
$\queryset.\emph{add}(\query)$
\label{slab:alg:toplevelalg:nocounterexample}
\Comment{$\emph{null}$ means $\query$ is equivalent to $\inputquery$.}

\Else \ $\counterexamples.\emph{add}( \counterexample )$;
\Comment{Otherwise, $\counterexample$ is a new counterexample.}
\label{slab:alg:toplevelalg:addcounterexample}

\EndIf

\EndFor

\vspace{-2pt}
\State $\cost \assign \cost - 1$;
\label{slab:alg:toplevelalg:decreasescore} 

\EndWhile

\vspace{-1pt}
\State $\query \assign \textsc{RankPerformance}(\queryset, \inputquery)$;
\label{slab:alg:toplevelalg:checkfaster}
\Comment{{{\scriptsize \Circled{3}}} Performance Ranker.}

\State \Return $\query$;
\label{slab:alg:toplevelalg:returnquery}

\caption{Top-level algorithm of \toolname.}
\label{slab:alg:toplevelalg}
\end{algorithmic}
\end{algorithm}
\vspace{-20pt}
\end{figure}

\end{comment}



\begin{example}
Suppose $\inputquery$ is query $\query_1$ from Table~\ref{slab:tab:Q1-Q2}. 
Our CEGIS-based algorithm will begin with query candidates with score 0: $\query'$ below is one such query. 
$\query'$ is obviously not equivalent to $\inputquery$, and our equivalence checker gives a counterexample $\counterexample$ below. Note that $\counterexample$ consists of an input DB $\inputdb$ and the output that $\inputquery$ returns on $\inputdb$. 
A query eventually returned by our tool will pass this counterexample. 



\begin{table}[H]
\footnotesize
\caption*{A counterexample $\counterexample$ consists of an input database $\inputdb$ (left) and an output table of $\inputquery$ on $\inputdb$.}
\vspace{-7pt}
\hspace*{-5pt}
\begin{tabular}{cc}
\makecell[l]{ $\query': $ \ \sqlselectcolor gender, day, \\  
\hspace{42pt} \sqlsumcolor(score) 
  \\ \hspace{2em} \sqlfromcolor scores
  \\ \hspace{2em} $\sqlgroupbycolor$ gender, day}
  &
\makecell[l]{ $\query\doubleprime: $ \sqlselectcolor gender, day, \\
\hspace{5.3em} \sqlsumcolor(score) \\  \hspace{5.3em} \sqlovercolor ($\sqlpartitionbycolor$ gender \\ \hspace{7.9em} $\sqlorderbycolor$ day)\\
 \hspace{2em}  \  \sqlfromcolor scores \\ \hspace{2em} \  $\sqlgroupbycolor$ gender, day}
\end{tabular}
\vspace{-12pt}
\end{table}



\begin{table}[H]
\footnotesize
\begin{tabular}{cc}
\begin{subtable}{.4\linewidth}
\begin{tabular}{|c|c|c|}
 \hline
        gender & day & score \\
        \hline
        a & 1 & 1 \\
        a & 2 & 2 \\
        \hline
\end{tabular}
\end{subtable}
&
\begin{subtable}{.4\linewidth}
\begin{tabular}{|c|c|c|}
 \hline
        gender & day & score \\
        \hline
        a & 1 & 1 \\
        a & 2 & 3 \\
        \hline
\end{tabular}
\end{subtable}
\end{tabular}
%\caption{counterexample $\counterexample$}

\vspace{-5pt}
\label{slab:tab:alg:counterexample}
\end{table}



On the other hand, the query $\query\doubleprime$ next to $\query$ above is equivalent to $\inputquery$. However, $\query\doubleprime$ is determined to be slower than $\query_2$ from Table~\ref{slab:tab:Q1-Q2} by our performance ranker. 
$\toolname$ was able to find both $\query_2$ and $\query\doubleprime$ which are added to $\queryset$; however, \toolname returns $\query_2$ in the end, since it is ranked the highest. 


\label{slab:ex:cegis}
\end{example}


\subsection{Dataflow-Based Query Score Function}
\label{slab:sec:scorefunction}

As we can see, a key challenge underlying the success of Algorithm~\ref{slab:alg:toplevelalg} is how to design a good score function as well as how to develop an algorithm that can effectively use the score function to prioritize the search. 
This section first addresses the score function design. 


\head{Key idea: scoring queries based on dataflows}
Recall that we use a score function to quantitatively rank queries in terms of their likelihood of optimizing a user-provided input query $\inputquery$. That is, given $\inputquery$ and $\query$, a desired score function should assign $\query$ with a higher score if $\query$ is indeed semantically equivalent to and faster than $\inputquery$. At the same time, it should minimize false positives: that is, it should not assign high scores to too many inequivalent queries. 


To design such a score function, our key insight is that, a query $\query$ that indeed optimizes $\inputquery$ typically exhibits \emph{dataflows} that are also manifested in $\inputquery$. 
For example, if $\query$ filters an aggregated column (e.g., using $\sqlhaving \sqlspace \sqlsum$), the same computations (e.g., calculating $\sqlsum$) likely would also show up in $\inputquery$, although the computations might be organized \emph{syntactically differently} (e.g., first calculating $\sqlsum$ in $\sqlselect$ and then using $\sqlwhere$ to filter, which is slower). 
In other words, our dataflow-based score function is designed to favor queries $\query$ that involve dataflows from $\inputquery$ --- the rationale is that, since $\inputquery$ is functionally correct, its underlying computations are likely sufficient for generating an optimized query.


However, we do not require $\query$ to use up all dataflows from $\inputquery$. In other words, $\inputquery$ may exhibit redundant dataflows that are not necessary for computing the same result and thus can be optimized away. For example, if $\inputquery$ involves a redundant $\sqljoin$, an optimized query $\query$ may contain only dataflows from one branch of the $\sqljoin$. 


Finally, if a query $\query'$ contains dataflows that are not exhibited in $\inputquery$, oftentimes $\query'$ is not equivalent and hence we can assign it with a lower score. For example, suppose $\inputquery$ calculates $\sqlsum$ over a column. It is less likely that a query $\query'$ can accomplish the same goal with only $\sqlavg$. Similarly, if $\inputquery$ and $\query'$ both perform filtering but on two different columns that are unrelated, it is more plausible to believe they are not semantically equivalent. 
%However, special care needs to be taken for ``alternative dataflows'' --- that is, two syntactically different dataflows may share the same underlying logic. 
%For instance, $\sqlpartitionby$ and $\sqlgroupby$, while syntactically different, can both be used for grouping. 



\barzanrev{Note that, 
since our score is calculated based on semantic information (i.e., dataflows) instead of syntactic features, a syntactically simpler query (e.g., one with fewer lines) may receive a lower score than a more complex query. 
In other words, our algorithm first prioritizes generating queries with higher scores, and 
then prioritizes syntactically smaller queries when their scores are the same (see Section~\ref{slab:sec:prioritizingsearch}). This is an advantage since it will allow us to uncover non-trivial rewrites sooner if they are more promising.} 


\head{Dataflows}
Our observations suggest a score function design that takes two queries (i.e., $\inputquery$ to be optimized and $\query$ to be scored), extracts dataflows for each, and calculates the score based on these dataflows. 
In what follows, we first formalize the notion of dataflow.
Then, we present the dataflow extraction algorithm. 

\begin{definition}[Dataflow]
A dataflow $\dataflowchain$ (defined in Figure~\ref{slab:fig:dataflowlanguage}) for a given query $\query$ is a sequence of SQL operations that are performed during $\query$'s execution on one or more input tables or their columns.
\end{definition}


Let us explain our dataflow language more formally (see Figure~\ref{slab:fig:dataflowlanguage}).
In the base cases, a dataflow $\dataflowchain$ is either an input table $\dbtable$ or a column from an input table. 
%it must not be a column alias created for an intermediate table. 
%The reason is to make dataflows independent on specific intermediate results of a query. 
In the recursive cases, $\dataflowchain$ captures operations performed on top of such base dataflows. 
For example, $\agg(\dataflowchain, \mydots, \dataflowchain)$ describes the application of an aggregate function over dataflows for argument columns. 
$[\dataflowchain, \mydots, \dataflowchain]$ is the composition of columns to form a table (e.g., via $\sqlselect$). 
Expressions, such as $=$ and $\leq$, involves comparison logics, which is what $\emph{op}(\dataflowchain, \dataflowchain)$ is designed for. The language also includes boolean combinations of dataflows in order to represent dataflows from the filtering conditions in a query.
Finally, we capture other SQL operations, such as $\sqlpartitionby$, $\sqlorderby$, $\sqlgroupby$, etc.
%Note that we do not include all SQL operators (e.g., $\sqlselect$) since $\sqlselect$ is already implicitly taken into account by $[\dataflowchain, \mydots, \dataflowchain]$. 

\begin{example}
In this example, we show some sample dataflows for $\inputquery$ from Example~\ref{slab:ex:cegis}:
$$
\small
 \left \{
\text{\texttt{scores.score}, $\texttt{scores.day} \leq \texttt{scores.day}$, \sqlsumcolor(\texttt{scores.score}), $\dotsb$ }
\right \}
$$
We have \texttt{scores.score} because $\query_{\emph{in}}$ uses data from \texttt{score} column in \texttt{scores} table. 
We have \sqlsumcolor(\texttt{scores.score}), since $\query_{\emph{in}}$ performs summation on this column. 
Similarly, $\texttt{scores.day} \leq \texttt{scores.day}$ because an intermediate step of $\query_{\emph{in}}$ does a comparison between these two columns. 
%As we can see, $\query_{\emph{in}}$'s dataflows capture all operations in $\query_{\emph{in}}$.

\begin{comment}
\begin{table}[H]
\small
\begin{tabular}{|c|c|}
 \hline
 $\dataflowset_{\emph{in}}$

 & 
 

   \makecell[l]{    % \sqldistinctcolor: 1 \ \ \ \ \ \ \  \hspace{1em}
  %                  \texttt{scores}: 2 \\
  %                  \texttt{scores.day}: 4 \ \ \ \ \hspace{1em}
  %                  \texttt{scores.gender}: 4\\
                   \texttt{scores.score}, \hspace{1em}
                   \sqlsumcolor(\texttt{scores.score})\\
                   % \texttt{scores.gender = scores.gender}: 1\\
                   \texttt{scores.day <= scores.day},\\
                   \sqlgroupbycolor \texttt{scores.gender},\\
                   \dotsb
                   % \sqlgroupbycolor \texttt{scores.day}: 1\\
              }  
\\
% \hline
% $\dataflowset_{2}$

%  & 
 
%  \makecell[l]{
%                 % \texttt{scores}: 1 \ \ \ \ \ \ \ \ \ \ \ \ \hspace{1em}
%                 % \texttt{scores.gender}: 2\\
%                 % \texttt{scores.day}: 2 \ \ \ \ \hspace{1em}
%                 % \texttt{scores.score}: 1\\
%                 \texttt{scores.score}, \ \ \ \ \hspace{1em}
%                 \sqlsumcolor(\texttt{scores.score}),\\
%                 \sqlpartitionbycolor \texttt{scores.gender},\\
%                 \sqlorderbycolor \texttt{scores.day}\\
%                 \dotsb
%               }  
% \\

\hline

\end{tabular}
\caption{Dataflows for $\query_{in}$ \xinyu{this can be put in one line.}}
\label{slab:table:runningexample:dataflow}
\end{table}
\end{comment}
% \todo{the goal here is to illustrate the dataflow idea. ideally we use motivating example. e.g., we can take the output query from the motivating example, show a few of its dataflows. note: do not explain how to extract the dataflows, but just show the dataflows. we can also show a couple of dataflows for the input query. @Jie/Rui?}
\end{example}

\begin{comment}


==== 


For instance, \rui{ $\leq$ and \sqlorderbycolor have the same semantics of ordering, \sqlpartitionbycolor and \sqlgroupbycolor have the same semantics of grouping. They can sometimes be used interchangeably, as shown in examples shown in Table \ref{slab:tab:Q1-Q2}.}

Overall, these observations suggest a design of the score function that takes two queries (i.e., $\inputquery$ to be optimized and $\query$ to be scored), extracts dataflows for each, and calculates the score based on these dataflows. 
In what follows, we first formalize the notion of dataflow.
Then, we present the dataflow extraction rules. 
Finally, we describe how to compute the score from dataflows. 

\para{Dataflow} 
A dataflow $\dataflowchain$ is defined as follows. 

\barzan{this def can be made a bit clearer by defining sub-program}
\begin{definition}[Dataflow]
A dataflow $\dataflowchain$ exhibited in a query $\query$ is a sequence of operations performed by some sub-program of $\query$ on one or more input tables or their columns.
\end{definition}

To be more concrete, our dataflow language is given in Figure~\ref{slab:fig:dataflowlanguage}.
In the base cases, a dataflow $\dataflowchain$ is either an input table $\dbtable$ or a column from $\dbtable$; it must not be a column alias created for an intermediate table. 
The reason is to make dataflows independent on specific intermediate results of a query. 
Furthermore, $\dataflowchain$ also captures operations performed on top of such base dataflows. 
For example, $\agg(\dataflowchain, \mydots, \dataflowchain)$ describes the application of an aggregate function over dataflows for argument columns. 
$[\dataflowchain, \mydots, \dataflowchain]$ characterizes the composition of columns to form a table (e.g., through $\sqlselect$). 
Expressions, such as $=$ and $\leq$, involves comparison logics, which is what $\emph{op}(\dataflowchain, \dataflowchain)$ is designed for. The language also includes boolean combinations of dataflows in order to represent dataflows from the filtering conditions in a query.
Finally, a dataflow also captures SQL operations, including $\sqlpartitionby$, $\sqlorderby$ and $\sqlgroupby$. 
Note that we do not include all SQL operators (e.g., $\sqlselect$) since $\sqlselect$ is already implicitly taken into account by $[\dataflowchain, \mydots, \dataflowchain]$. 


\begin{example}
\rui{please review here} Table \ref{slab:table:runningexample:dataflow} belows shows dataflows $\dataflowset_{1}$, $\dataflowset_{2}$ extracted from \query1 and \query2 from Table \ref{slab:tab:Q1-Q2}. Dataflows are represented in a multi-set. where each dataflow can appear multiple times. We will show how to extract them in details later. Take $\dataflowset_{1}$ as an example, \sqlsumcolor(\texttt{scores.score}) means a summation on \texttt{score} column of \texttt{scores} table is performed by $\query1$, \texttt{scores.day <= scores.day} means a comparison between \texttt{day} columns of \texttt{scores} table is done in $\query1$. These are clues that capture the semantics of a query.


\begin{table}[H]
\small
\begin{tabular}{|c|c|}
 \hline
 $\dataflowset_{1}$

 & 
 

   \makecell[l]{    % \sqldistinctcolor: 1 \ \ \ \ \ \ \  \hspace{1em}
  %                  \texttt{scores}: 2 \\
  %                  \texttt{scores.day}: 4 \ \ \ \ \hspace{1em}
  %                  \texttt{scores.gender}: 4\\
                   \texttt{scores.score}, \hspace{1em}
                   \sqlsumcolor(\texttt{scores.score})\\
                   % \texttt{scores.gender = scores.gender}: 1\\
                   \texttt{scores.day <= scores.day},\\
                   \sqlgroupbycolor \texttt{scores.gender},\\
                   \dotsb
                   % \sqlgroupbycolor \texttt{scores.day}: 1\\
              }  
\\
\hline
$\dataflowset_{2}$

 & 
 
 \makecell[l]{
                % \texttt{scores}: 1 \ \ \ \ \ \ \ \ \ \ \ \ \hspace{1em}
                % \texttt{scores.gender}: 2\\
                % \texttt{scores.day}: 2 \ \ \ \ \hspace{1em}
                % \texttt{scores.score}: 1\\
                \texttt{scores.score}, \ \ \ \ \hspace{1em}
                \sqlsumcolor(\texttt{scores.score}),\\
                \sqlpartitionbycolor \texttt{scores.gender},\\
                \sqlorderbycolor \texttt{scores.day}\\
                \dotsb
              }  
\\

\hline

\end{tabular}
\caption{Dataflows for \query1 and \query2}
\label{slab:table:runningexample:dataflow}
\end{table}
\end{example}

% \todo{the goal here is to illustrate the dataflow idea. ideally we use motivating example. e.g., we can take the output query from the motivating example, show a few of its dataflows. note: do not explain how to extract the dataflows, but just show the dataflows. we can also show a couple of dataflows for the input query. @Jie/Rui?}
\end{comment}




\begin{figure}[!t]
\vspace{10pt}
\footnotesize 
\centering
\[
\arraycolsep=1.5pt\def\arraystretch{1.1}
\begin{array}{rrll}
&
\dataflowchain & 
\grammareq & 
\emph{an input table} \ | \ 
\emph{a column from an input table}  \\ 
& & \ \ \ | &  
\agg(\dataflowchain, \mydots, \dataflowchain) \ | \ 
\windowfunc \ | \ 
[\dataflowchain, \mydots, \dataflowchain] \ | \ 
\emph{op}(\dataflowchain, \dataflowchain) \ | \ 
\dataflowchain \land \dataflowchain \ | \ 
\dataflowchain \lor \dataflowchain \\ 
& & \ \ \ | &  
\sqlpartitionby(\dataflowchain) \ | \ 
\sqlorderby(\dataflowchain) \ | \ 
\sqlgroupby(\dataflowchain) \\ 
\end{array}
\]
\vspace{-12pt}
\caption{Dataflow language.}
\label{slab:fig:dataflowlanguage}
\vspace{-12pt}
\end{figure}





\begin{figure*}[!t]
\footnotesize 
\resetmycounter{mycounter}
% NOTE (dissertation): content was sized for acmart's full page width; scaled to fit.
\begin{center}\resizebox{\textwidth}{!}{$\displaystyle
\arraycolsep=.5pt\def\arraystretch{1.1}
\begin{array}{l}
% first block 
\begin{array}{lll}
\fbox{\textsc{AllDfs} \text{algorithm that extracts all dataflows for queries}}
\\[5pt]
\begin{array}{cc}
(\showmycounter{mycounter}) & 
\irule{
\dbtable \text{ is an input table}
}{
\calcalldataflowchains{}{(\dbtable)}{ \{ \dbtable \} }
}
\end{array}
\begin{array}{cc}
(\showmycounter{mycounter}) & 
\irule{
\calcalldataflowchains{}{(\query)}{\dataflowset_1}  \ \ \ \ 
\calcalldataflowchains{\query}{(\sqlpredicate)}{\dataflowset_2} \ \ \ \ 
\calcalldataflowchains{\query}{(\targetlist)}{\dataflowset_3}
}{
\calcalldataflowchains{}{(\sqlselect \sqlspace \targetlist \sqlspace \sqlfrom \sqlspace \query \sqlspace \sqlwhere \sqlspace \sqlpredicate)}{\dataflowset_1 \cup \dataflowset_2 \cup \dataflowset_3}
}
\end{array}
\begin{array}{cc}
(\showmycounter{mycounter}) & 
\irule{
\calcalldataflowchains{}{(\query_1)}{\dataflowset_1} \ \ \ \ 
\calcalldataflowchains{}{(\query_2)}{\dataflowset_2} \ \ \ \ 
\calcalldataflowchains{\query_1 \sqlspace \sqljoin \sqlspace \query_2}{(\sqlpredicate)}{\dataflowset_3}
}{
\calcalldataflowchains{}{(\query_1 \sqlspace \sqljoin \sqlspace \query_2 \sqlspace \sqlon \sqlspace \sqlpredicate)}{ \dataflowset_1 \cup \dataflowset_2 \cup \dataflowset_3 } 
}
\end{array}
\end{array}
\\ \\ 
% second block 
\begin{array}{l}
\begin{array}{lll}
\fbox{\textsc{AllDfs} \text{algorithm that extracts all dataflows for conditions}}
\\[5pt]
\begin{array}{cc}
(\showmycounter{mycounter}) & 
\irule{
\calcalldataflowchains{\query}{(\sqlpredicate_i)}{\dataflowset_i} \ \ \ \ 
\calcdataflowchain{\query}{(\sqlpredicate_i)}{\dataflowchain_i} 
}{
\calcalldataflowchains{\query}{(\sqlpredicate_1 \ \emph{lop} \ \sqlpredicate_2)}{\dataflowset_1 \cup \dataflowset_2 \cup \{\emph{lop}(\dataflowchain_1, \dataflowchain_2) \}}
}
\end{array}
\begin{array}{cc}
(\showmycounter{mycounter}) & 
\irule{
\calcalldataflowchains{\query}{(\sqlexpr_i)}{\dataflowset_i} \ \ \ \ 
\calcdataflowchain{\query}{(\sqlexpr_i)}{\dataflowchain_i} 
}{
\calcalldataflowchains{\query}{(\sqlexpr_1 \emph{ op } \sqlexpr_2)}{\dataflowset_1 \cup \dataflowset_2 \cup \{ \emph{op}(\dataflowchain_1 , \dataflowchain_2) \}}
}
\end{array}
\\ \\ 
\fbox{\textsc{Df} \text{algorithm that extracts dataflows for conditions, expressions and targets}}
\\[5pt]
\begin{array}{cc}
(\showmycounter{mycounter}) & 
\irule{
\calcdataflowchain{\query}{(\sqlpredicate_i)}{\dataflowchain_i}
}{
\calcdataflowchain{\query}{(\sqlpredicate_1 \ \emph{lop} \  \sqlpredicate_2)}{ \emph{lop}(\dataflowchain_1, \dataflowchain_2) }
}
\end{array}
\begin{array}{cc}
(\showmycounter{mycounter}) & 
\irule{
\calcdataflowchain{\query}{(\sqlexpr_i)}{\dataflowchain_i}
}{
\calcdataflowchain{\query}{(\sqlexpr_1 \emph{ op } \sqlexpr_2)}{ \emph{op}(\dataflowchain_1 , \dataflowchain_2)  }
}
\end{array}
\begin{array}{cc}
\vspace{-5pt}
(\showmycounter{mycounter}) & 
\irule{
}{
\calcdataflowchain{\query}{(\emph{const})}{\emph{const}}
}
\end{array}
\begin{array}{cc}
(\showmycounter{mycounter}) & 
\irule{
\dbtable \text{ is an input table}
}{
\calcdataflowchain{\dbtable}{(\col)}{{T.}\col}
}
\end{array}
\begin{array}{cc}
(\showmycounter{mycounter}) & 
\irule{
\col \in \textsc{Columns}(\query_i) \ \ \ \ 
%\col \not\in \textsc{Columns}(\query_{2-i}) \ \ \ \ 
\calcdataflowchain{\query_i}{(\col)}{\dataflowchain_i}
}{
\calcdataflowchain{\query_1 \sqlspace \sqljoin \sqlspace \query_2 \sqlspace \sqlon \sqlspace \sqlpredicate}{(\col)}{ \dataflowchain_i } 
}
\end{array}
\\ \\ 
\begin{array}{cc}
(\showmycounter{mycounter}) & 
\irule{
(\target \sqlspace \sqlas \sqlspace \col) \in \targetlist \ \ \ \ 
\calcdataflowchain{\query}{(\target)}{\dataflowchain}
}{
\calcdataflowchain{\sqlselect \sqlspace \targetlist \sqlspace \sqlfrom \sqlspace \query \sqlspace \sqlwhere \sqlspace \sqlpredicate}{(\col)}{ \dataflowchain } 
}
\end{array}
\begin{array}{cc}
(\showmycounter{mycounter}) & 
\irule{
\calcdataflowchain{\query}{(\col)}{\dataflowchain}
}{
\calcdataflowchain{\query}{(\agg(\col))}{\agg(\dataflowchain)}
}
\end{array}
\begin{array}{cc}
(\showmycounter{mycounter}) & 
\irule{
\calcdataflowchain{\query}{(\col)}{\dataflowchain} 
}{
\calcdataflowchain{\query}{\big( \agg(\col) \sqlspace \sqlover( \sqlpartitionby \sqlspace \cols_1 \sqlspace \sqlorderby \sqlspace \cols_2 ) \big)}{ \agg(\dataflowchain) } 
}
\end{array}

\end{array}
\end{array}
\end{array}
$}\end{center}
\vspace{-10pt}
\caption{Inference rules that explain how dataflow extraction works for some key constructs in our query language.}
\label{slab:fig:dataflowextraction}
\vspace{-5pt}
\end{figure*}


{\head{Extracting dataflows from queries}
Now let us talk about how to extract a set $\dataflowset$ of dataflows from a query $\query$. 
Figure~\ref{slab:fig:dataflowextraction} shows how to extract dataflows where we use a judgment of the form:
\[\centering
\calcalldataflowchains{\query_{\emph{ctx}}}{(\program)}{ \dataflowset }
\]
This means: given a ``context'' $\query_{\emph{ctx}}$ associated with $\program$, $\dataflowset$ is the set of \emph{all} dataflows that are manifested in $\program$. 
Here, $\program$ may be a query, a condition, a target list, etc. 
In general, $\program$ is associated with a context $\query_{\emph{ctx}}$ --- e.g., if $\program$ is a target list, $\query_{\emph{ctx}}$ is the query that produces the table that the targets in $\program$ correspond to. 
Having a context allows us to trace the flow of data to the original input tables.}

%Due to the space limit, Figure~\ref{slab:fig:dataflowextraction} shows how to extract dataflows for a key subset of our language. 
Let us first explain how to extract dataflows for a query $\query$. 
The first rule in Figure~\ref{slab:fig:dataflowextraction} concerns the base case where $\query$ is an input table $\dbtable$: in this case, the result is a singleton set with $\dbtable$. Intuitively, this means, the computation in $\query$ involves only $\dbtable$ and nothing else. 
The next one concerns selection --- it states that, the dataflow set $\dataflowset$ of a $\sqlselect$ query is the union of three sets: $\dataflowset_1$ for the query $\query$ to select from, $\dataflowset_2$ for the filtering condition $\sqlcondition$, and $\dataflowset_3$ for the target list $\targetlist$. Conceptually, $\dataflowset_1$ contains dataflows manifested in \emph{all components} of $\query$; that is, if $\query$ is a nested query, $\dataflowset_1$ would also include dataflows from the sub-queries. 
Similarly, $\dataflowset_2$ and $\dataflowset_3$ capture computations performed in $\sqlcondition$ and $\targetlist$. 
%Also note that, the selection query itself does not create new dataflows; its dataflows all come from its arguments $\query$, $\sqlcondition$ and $\targetlist$. 
(3) is very similar to (2) in that it also takes the union of dataflow sets for each of the arguments of $\sqljoin$. The join logic is implicitly captured in  $\dataflowset_3$ for the join condition $\sqlcondition$.

Next, let us examine the extraction rules for a condition $\sqlpredicate$, which can be used as a join condition or used in $\sqlwhere$ or $\sqlhaving$. 
Looking at (4), for a boolean combination $\emph{lop}$ (e.g., $\lor$) of multiple conditions, we would first recursively invoke $\textsc{AllDfs}$ to obtain \emph{all dataflows} $\dataflowset_i$ for each $\sqlcondition_i$. In addition, since $\emph{lop}$ itself performs a logical operation, the final dataflow set also includes $\emph{lop}(\dataflowchain_1, \dataflowchain_2)$: here, $\dataflowchain_i$ is the dataflow for $\sqlcondition_i$ that (different from $\dataflowset_i$) does \emph{not} include dataflows from $\sqlcondition_i$'s components, and we use a separate function $\textsc{Df}$ (different from $\textsc{AllDfs}$) to compute $\dataflowchain_i$. The key difference between $\textsc{AllDfs}$ and $\textsc{Df}$ is that, the former includes all dataflows for all components in $\program$, whereas the latter only concerns one dataflow for $\program$ itself. The next rule (5) takes care of the base case where the condition $\sqlcondition$ is an application of $\emph{op}$ (e.g., $=$) over expressions (e.g., columns). It also makes use of $\textsc{Df}$ to retrieve the dataflow for an expression $\sqlexpr_i$. 

Let us dive a little deeper to see how $\textsc{Df}$ works. 
(6) states that the dataflow of a logical operation (e.g., $\land$) is composed of dataflows of the two arguments $\sqlcondition_1, \sqlcondition_2$. (7) is very similar except that it concerns comparison operations (like $=$) and it invokes $\textsc{Df}$ on expressions $\sqlexpr_i$ such as a column, an aggregate function applied to a column, or a constant. 
(8)-(13) detail how to extract dataflows for expressions. 
(8) says the dataflow for a constant is the constant itself. 
(9) states that, given an \emph{input} table $\dbtable$, the dataflow for its column $\col$ is $\dbtable.\col$. 
(10) considers extracting the dataflow for a column $\col$ --- which may be a column alias --- given a $\sqljoin$ query: it recursively extracts the dataflow for $\col$ given $\query_i$ that $\col$ comes from. 
(11) is similar: it first identifies the target $\target$ from the target list $\targetlist$ that corresponds to $\col$ and then invokes $\textsc{Df}$ on $\target$ given the query $\query$ being selected from. 
Finally, (12) and (13) extract dataflows for aggregate functions. 

\begin{example}
Let us explain how to extract the dataflows for the following component (i.e., a comparison) in $\query_1$ from Table \ref{slab:tab:Q1-Q2}.
$$
\text{
\texttt{s2.day <= s1.day}}
$$

The dataflow set of this expression is the union of three sets of dataflows: (1) the dataflows of its left argument (\{\texttt{scores.day}\}), (2) the dataflows of the right argument (namely \{\texttt{scores.day}\}), and (3) the dataflow of the comparison (i.e., $\{\texttt{scores.day} \leq \texttt{scores.day}\}$). This essentially means that the query uses  data from \texttt{day} column of the input \texttt{scores} table and performs a comparison where the left and right arguments come from the \texttt{day} column of \texttt{scores} table as well. This provides clues to guide the search. 
%because equivalent and faster queries are very likely to preserve the same semantics.
\end{example}

\begin{comment}
\begin{example}
We show how to extract the dataflows of the subprogram 
$$
\small
\text{
\sqlsumcolor~(\texttt{score}) \sqlovercolor ( \sqlpartitionbycolor~\texttt{gender} \sqlorderbycolor~\texttt{day})}
$$
of $\query2$ in Table \ref{slab:tab:Q1-Q2}. According to (17), the dataflow of this subprogram is the union of dataflows of its three subprograms - \sqlsumcolor(\texttt{score}), \sqlpartitionbycolor \texttt{gender}, and \sqlorderbycolor \texttt{day}. Thus we need to recursively get the dataflows of its subprograms first. 

To get the dataflows of \sqlsumcolor(\texttt{score}), we apply rule (16) as follows. Note that we also need to apply rule (9) and rule (15) recursively, but we will skip here since it's straightforward:

$$
\footnotesize
\irule{
\calcalldataflowchains{\query}{(\text{\texttt{score}})}{\text{\{\texttt{scores.score}\}}} \ \ \ \ 
\calcdataflowchain{\query}{(\text{\{\texttt{score}\}})}{\text{\texttt{scores.score}}} 
}{
\calcalldataflowchains{\query}{\big( \text{\sqlsumcolor(\texttt{score})} \big)}{ {\text{\{\texttt{scores.score}\}}} \cup \{ \text{\sqlsumcolor(\texttt{scores.score})} \} } 
}
$$

Similarly, we can apply corresponding rules to get the dataflows of \sqlpartitionbycolor \texttt{gender} and \sqlorderbycolor \texttt{day}.


% \todo{this example should illustrate the extraction process for one candidate query in the motivating example. @Rui/Jie?}
\end{example}
\end{comment}


\begin{comment}

\head{Extracting alternative dataflows for input query}
\todo{@Rui/Jie: let's discuss how this works.}
\rui{I write starting here} As we discussed before, two syntactically different dataflows might share the same underlying logic. We call them \emph{alternative dataflows}.  Therefore, when we see some specific dataflows, we may also want to consider taking those alternative dataflows into account when calculating the scores. For example, when we see $\dataflowchain_1 \leq \dataflowchain_2$, we will treat $\sqlorderby \sqlspace \dataflowchain_1$ and $\sqlorderby \sqlspace \dataflowchain_2$ as the alternative chains since they preserve the same semantics of ordering. We will also treat $\sqlgroupby \sqlspace \dataflowchain$ and $\sqlpartitionby \sqlspace \dataflowchain$ as alternative chains interchangeably because they preserve the same semantics of grouping. \rui{do we show all alternative chains we have? Or just breifly discuss like this?}
\end{comment}


{\head{Dataflow-based scoring}
Now we define our $\scorefunction$ function, which scores a program $\program$ (potentially associated with a context $\query_{\emph{ctx}}$) with respect to an input query $\inputquery$. 
First, it extracts two sets of dataflows, i.e., $\dataflowset_{\emph{in}}$ and $\dataflowset$, for $\inputquery$ and $\program$ respectively using the algorithm from Figure~\ref{slab:fig:dataflowextraction}. 
Then, it computes the set difference, i.e., $\dataflowset_{\emph{diff}} = \dataflowset \setminus \dataflowset_{in}$, which gives the dataflows in $\program$ but not in $\inputquery$. 
Finally, we define $\scorefunction(\program | \query_{\emph{ctx}}, \inputquery) = -|\dataflowset_{\emph{diff}}|$. 
As we can see, $0$ is the highest possible score: in this case, all of $\program$'s dataflows are from $\inputquery$. 
For brevity, we use notation $\scorefunction(\program|\query_{\emph{ctx}})$ when $\inputquery$ is clear from the context.}




{\head{Monotonicity}
An important property of our $\scorefunction$ function is that it is \emph{monotonic}. 
That is, given $\inputquery$ and $\query_{\emph{ctx}}$, for any component $\program'$ of $\program$ (e.g., $\program'$ is a sub-query of query $\program$), $\scorefunction(\program'|\query_{\emph{ctx}}) \geq \scorefunction(\program|\query_{\emph{ctx}})$. 
The proof is obvious due to the monotonic nature of our dataflow extraction algorithm and our $\scorefunction$ function: the dataflow set for $\program'$ is always a subset of that for $\program$, hence $\program$'s score is no less than $\program$'s. 
The implication is: when composing multiple programs $\program_1, \mydots, \program_n$ to form a bigger program $\program$, $\program$'s score is never higher than the score of any $\program_i$'s. 
The next Section~\ref{slab:sec:prioritizingsearch} presents an algorithm that uses this property to effectively guide the query synthesis process.}









\subsection{Prioritizing Search using Score Function}
\label{slab:sec:prioritizingsearch}


%\revisioncancut{So far, \ref{slab:sec:scorefunction} has addressed the first challenge of how to score a query $\query$ with respect to a user-provided input query $\inputquery$. 
%In this section, we address the second challenge: how to leverage our score function to effectively guide the query synthesis process?}

%\head{Key idea: leveraging monotonicity to stratify search space}
\barzanrev{Our key idea is to leverage the \emph{monotonicity} property of our score function to develop a \emph{stratified search algorithm} that enumerates programs in layers: it first finds all queries with score $0$ (i.e., the highest possible score) \emph{without} considering those with lower scores, then generates all queries with score $-1$ (i.e., the second highest) \emph{again without} constructing those with lower scores, and so on.}



\begin{figure}[!t]
\footnotesize 
\centering

\vspace{-6pt}
\begin{algorithm}[H]
\begin{algorithmic}[1]

\myprocedure{LazyEnumerate}{$\inputquery$}

\Statex\Input{Input query $\inputquery$.}


\Statex\Output{A stream of candidate queries in non-ascending order of scores.}

\vspace{2pt}

\State 
$\queryset_{\geq 1} \assign \emptyset$;  
\label{slab:alg:lazyenumerate:init}
\Comment{All queries scored at least 1, which is empty.}


\For{$\score = 0, -1, -2, \mydots$} 
\label{slab:alg:lazyenumerate:forbegin}


\State
$\queryset_{\geq \score} \assign \queryset_{\geq \score+1}$;
\label{slab:alg:lazyenumerate:initS}
\Comment{All queries scored $\geq \score$, initially set to $\queryset_{\geq \score+1}$.}


\While{true}
\label{slab:alg:lazyenumerate:whilebegin}

\State
$\queryset'_{\score} \assign \textsc{NewQueriesWithScoreS}(\inputquery, \queryset_{\geq \score}, \score)$; 
\label{slab:alg:lazyenumerate:constructnew}

\If{$\queryset'_{\score} = \emptyset$} 
\textbf{break};
\Comment{Exit if no more new queries scored $\score$.}
\EndIf
\label{slab:alg:lazyenumerate:whilebreak}


\vspace{-1pt}
\State
\textbf{yield} $\queryset'_{\score}$; 
\label{slab:alg:lazyenumerate:yield}


\vspace{1pt}
\State
$\queryset_{\geq \score} \assign \queryset_{\geq \score} \cup \queryset'_{\score}$; 
\label{slab:alg:lazyenumerate:merge}
\Comment{All queries scored at least $\score$ so far.}

\EndWhile
\label{slab:alg:lazyenumerate:whileend}

\EndFor
\label{slab:alg:lazyenumerate:forend}


\caption{{$\textsc{LazyEnumerate}$ search algorithm.}}
\label{slab:alg:lazyenumerate}
\end{algorithmic}
\end{algorithm}
\vspace{-20pt}
\end{figure}




\barzanrev{In particular, to generate \emph{all} queries $\queryset_{\score}$ whose score is exactly $\score$, we track only those queries $\queryset_{\geq \score}$ whose score is \emph{at least} $\score$, because according to the monotonicity property,   $\queryset_{\geq \score}$ is  sufficient to construct $\queryset_{\score}$. 
This allows for a dynamic programming design that lazily enumerates queries, as presented in Algorithm~\ref{slab:alg:lazyenumerate}.
%ta inja
It takes the input query $\inputquery$ and returns a stream of queries in a non-ascending order of their scores. 
Line~\ref{slab:alg:lazyenumerate:init} initializes $\queryset_{\geq 1}$ to empty set since the highest possible score is 0. 
Then, the loop at lines~\ref{slab:alg:lazyenumerate:forbegin}--\ref{slab:alg:lazyenumerate:forend} populates each $\queryset_{\geq \score}$ for all possible scores in descending order from 0.
Specifically, given $\score$, $\queryset_{\geq\score}$ is initialized to $\queryset_{\geq\score+1}$ (line~\ref{slab:alg:lazyenumerate:initS}). 
The inner loop (lines~\ref{slab:alg:lazyenumerate:whilebegin}--\ref{slab:alg:lazyenumerate:whileend}) iteratively constructs new queries $\queryset'_{\score}$ with score exactly $\score$ from $\queryset_{\geq\score}$ (line~\ref{slab:alg:lazyenumerate:constructnew}) until no more new queries are generated (line~\ref{slab:alg:lazyenumerate:whilebreak}). 
Line~\ref{slab:alg:lazyenumerate:yield} yields new queries $\queryset'_{\score}$ which are then added to $\queryset_{\geq\score}$ at line~\ref{slab:alg:lazyenumerate:merge}. Note that the while loop in Algorithm~\ref{slab:alg:lazyenumerate} may be non-terminating if the score function is based on \emph{sets} of dataflows: it is possible to keep generating new queries with the same set of dataflows (hence the same score). To ensure termination, dataflows from a program are represented as a \emph{multiset} where duplicate dataflows are counted multiple times; therefore, one cannot keep constructing new queries without decreasing the score.}



\begin{figure}[!t]
%\small 
\small 
\resetmycounter{mycounter}
\begin{centermath}
\arraycolsep=.5pt\def\arraystretch{.8}
\begin{array}{l}
(\showmycounter{mycounter}) \ 
\irule{
\dbtable \text{ is input table} \ \ \ 
\scorefunction(\dbtable) = \score
}{
\dbtable \in \queryset'_{\score}
}
\ \ 
(\showmycounter{mycounter}) \ 
\irule{
\query_i \in \queryset_{\geq \score} \ \ \ 
\sqlcondition \in \sqlconditionset_{\geq \cost}(\query) \ \ \ 
\scorefunction(\query') = \score
}{
( \query' = \query_1 \sqlspace \sqljoin \sqlspace \query_2 \sqlspace \sqlon \sqlspace \sqlcondition ) \in \queryset'_{\geq \score}
}
\\ \\ 
% (\showmycounter{mycounter}) \ 
% \irule{
% \sqlcondition_i \in \sqlconditionset_{\geq \score}(\query)
% }{
% \sqlcondition_1 \emph{ lop } \sqlcondition_2 \in \sqlconditionset_{\geq \score}(\query)
% }

(\showmycounter{mycounter}) \ 
\irule{
\sqlcondition_i \in \sqlconditionset_{\geq \score}(\query) \ \ \ \ \scorefunction(\sqlcondition' | \query) \geq \score
}{
( \sqlcondition' = \sqlcondition_1 \emph{ lop } \sqlcondition_2 ) \in \sqlconditionset_{\geq \score}(\query)
}
\ \ 
% (\showmycounter{mycounter}) \ 
% \irule{
% \sqlexpr_i \in \sqlexprset_{\geq \score}(\query)
% }{
% \sqlexpr_1 \emph{ op } \sqlexpr_2 \in \sqlconditionset_{\geq \score}(\query)
% }
(\showmycounter{mycounter}) \ 
\irule{
\sqlexpr_i \in \sqlexprset_{\geq \score}(\query) \ \ \ \ \scorefunction(\sqlcondition' | \query) \geq \score
}{
( \sqlcondition' = \sqlexpr_1 \emph{ op } \sqlexpr_2 ) \in \sqlconditionset_{\geq \score}(\query)
}
\\ \\ 
(\showmycounter{mycounter}) \ 
\irule{
\emph{const} \text{ is used in } \inputquery
}{
\emph{const} \in \sqlexprset_{\geq \score}(\query)
}
\ \ 
(\showmycounter{mycounter}) \ 
\irule{
\col \in \emph{Cols}(\query) \ \ \ \ 
\scorefunction(\col | \query) \geq \score
}{
\col \in \sqlexprset_{\geq \score}(\query)
}
\\ \\ 
(\showmycounter{mycounter}) \ 
\irule{
\query \in \queryset_{\geq \cost} \ \ \ \ 
\targetlist \in \targetlistset_{\geq \cost}(\query) \ \ \ \ 
\sqlcondition \in \sqlconditionset_{\geq \cost}(\query) \ \ \ \ 
\scorefunction(\query') = \score
}{
( \query' = \sqlselect \sqlspace \targetlist \sqlspace \sqlfrom \sqlspace \query \sqlspace \sqlwhere \sqlspace \sqlpredicate ) 
\in \queryset'_{\score}
}
\\ \\ 
(\showmycounter{mycounter}) \ 
\irule{
\target_i \in \targetset_{\geq \score}(\query) \ \ \ \ 
\colalias_i \text{ is a fresh column alias} \ \ \ \ 
\scorefunction(\targetlist | \query) \geq \score
}{
( \targetlist = [ \target_1 \sqlspace \sqlas \sqlspace \colalias_1, \mydots, \target_n \sqlspace \sqlas \sqlspace \colalias_n ] ) \in \targetlistset_{\geq \score}(\query)
}
\end{array}
\end{centermath}
\vspace{-5pt}
\caption{{Inference rules for \textsc{NewQueriesWithScoreS}.}
}
\vspace{-10pt}
\label{slab:fig:newqueries}
\end{figure}






\barzanrev{Next, let us explain how \textsc{NewQueriesWithScoreS} (invoked at line~\ref{slab:alg:lazyenumerate:constructnew} of Algorithm~\ref{slab:alg:lazyenumerate}) works. 
Due to space limit, Figure~\ref{slab:fig:newqueries} shows a key subset% NOTE (dissertation merge): the exported Overleaf copy contained a post-camera-ready footnote here (`\footnote{\barnzarev{For the complete set of rules, please see~\cite{tr-slabcity}.}}`) using an undefined macro \barnzarev; it does not appear in the published PVLDB version, so it is omitted. TODO: in Phase 2, consider including the complete rule set (no space limits in a dissertation).
 of inference rules that describe how to construct a new query $\query'$ with score $\score$ from an existing set $\queryset_{\geq\score}$ of queries with at least score $\score$.
Rule 1 is the base case stating that an input table $\dbtable$ is generated if $\dbtable$ has score $\score$. Rule 2 describes a recursive case: we join $\query_1$ and $\query_2$ using condition $\sqlcondition$ 
to generate a $\sqljoin$ query $\query'$ if 
$\query_1$ and $\query_2$
are both in $\queryset_{\geq\score}$,  
$\sqlcondition$ has a score of at least $\score$ given context $\query$, and the score of $\query'$ is exactly $\score$. 
Rules 3--6 explain how conditions with a score of at least $\score$ can be generated, where
Rules 5--6 provide the base cases and Rules 3--4 recursively build larger conditions. 
Rule 7 is an example to illustrate how to generate $\sqlselect$ queries, which additionally involves generating target lists $\targetlist$ with score at least $\score$.
Rule 8 explains how to generate such target lists. 
Note that, while implicit in the rule specification, we return $\query'$ only if it is not already in $\queryset_{\geq\score}$.}
% \barzan{changed all refs to numbered rules to Rules. lets make sure we do it everywhere consistently}

\barzanrev{In summary, new queries are generated in a bottom-up manner. That is, when synthesizing queries with score $\score$, we only use queries with a score of at least $\score$ as building blocks. Further, among queries with the same score, 
we prioritize syntactically smaller queries.}


 

\subsection{Checking Query Equivalence}
\label{slab:sec:checkingequiv}

So far, we have seen how the inductive synthesizer in the CEGIS-based algorithm works. In this section, we explain how the equivalence checker works. 
As mentioned earlier, \slabcity incorporates a range of SQL query equivalence checking techniques, including empirical testing as well as formal verification. 
These techniques are complementary to each other: testing is relatively cheap but cannot prove equivalence, while verification can provide an equivalence guarantee but tends to be significantly more expensive. 

\head{Syntax-based query testing}
To the best of our knowledge, there is little work on query equivalence testing for an expressive SQL language like ours. Thus, we develop a new tester specifically targeting our complex language. 
%with the goal of identifying counterexamples efficiently for as many incorrect queries as possible. 
Given an input query $\inputquery$ and a candidate query $\query$, as well as an integrity constraint $\dataconstraint$, the tester returns either (1) a database $\inputdb$ which satisfies $\dataconstraint$ such that $\inputquery$ and $\query$ produce different tables on $\inputdb$ or (2) ``unknown'' indicating it is not able to find a counterexample. 
In the latter case, it is possible that $\inputquery$ and $\query$ are indeed equivalent, or they are not equivalent but the tester is not able to disprove it. The key challenge is how to  disprove many non-equivalent query pairs from an expressive language without being too costly. 
Randomly generating test inputs does not work.
%due to the sparsity of the input space. 

Our key observation is that, we can extract hints from $\inputquery$ and $\query$ using a mostly syntactic analysis to effectively guide the input generation process. Below are some of our key insights. 
\begin{itemize}[leftmargin=*]
\item 
\emph{Leveraging filter conditions.}
We observe that databases producing \emph{non-empty} outputs tend to be useful for disproving equivalence. 
Therefore, we extract filter conditions (such as those in $\sqlwhere$) from $\inputquery$ and populate the test database with rows that satisfy these conditions. 
We also collect the $\sqljoin$ conditions to generate inputs that do not produce empty intermediate join results. 
%A simple yet effective heuristic we use is to enforce the $\sqljoin$ mapping between tables. Specifically, we require that (1) the tables to be joined have the same number of rows and (2) the $i$-th row of one table is joined with the $i$-th row of the other table.
%For self-joins, we split the table into two halves, which reduces the problem to that of a normal join. 
\item 
\emph{Duplicating non-key values.}
A second source of non-equivalence that we have observed is from wrong logic around operations like addition/removal and from certain keywords such as \sqldistinct being placed in a wrong location. 
Therefore, our tester identifies such operations and keywords to generate potentially useful test inputs accordingly. For example, if $\query$ has \sqldistinct on a column $C$, we will generate a test DB with duplicated values in $C$, with the goal of triggering the $\sqldistinct$ logic. 
\item 
\emph{Duplicating $\sqlgroupby$ columns.}
Similar to the previous insight, we also observed a common class of incorrect query candidates with wrong \sqlgroupby clauses. 
While certain grouping strategies definitely lead to equivalent queries that are significantly faster, some of them produce non-equivalent ones. 
Therefore, our tester looks for columns in \sqlgroupby of $\query$ that are different from those in $\inputquery$. Then, it generates test DBs that contain duplicated values in such columns. 
%We found these test inputs useful because duplicated values from the grouped columns will eventually be eliminated, which means such inputs are likely to produce different outputs for two queries with different \sqlgroupby clauses. 
\end{itemize}

In general, each insight above leads to a set of test input DBs, all of which are encoded into a logical formula $\varphi_1$. We also encode the integrity constraint into a formula $\varphi_2$. 
We use a constraint solver (such as Microsoft Z3) to generate satisfying assignments for $\varphi_1 \land \varphi_2$, where each satisfying assignment corresponds to one DB instance. 
\barzanrev{If $\varphi_1$ and $\varphi_2$ are contradictory (i.e., $\varphi_1 \land \varphi_2$ is unsatisfiable), it means no valid test inputs can be generated to differentiate $\inputquery$ and $\query$. In that case, we resort to subsequent verification approaches.}







\head{Bounded verification}
In addition to testing, we also use a bounded verifier to check equivalence and find counterexamples: it considers all inputs in a bounded space and thus is more exhaustive than our tester. 
\toolname incorporates existing bounded verifiers (i.e., those in \cosette~\cite{chu2017cosette,wang2018speeding}). 
In particular, if verified, it guarantees that $\inputquery$ and $\query$ produce the same output for \emph{all} input databases that have up to $k$ rows ($k$ is set to 2 in \toolname and can be customized by users), though cell values are not bounded. 
These queries will then be sent to our full verifier which we explain next. 


\head{Full verification}
\toolname incorporates state-of-the-art 
\emph{full-fledged verifiers} from \cosette~\cite{chu2017cosette} and \spes~\cite{zhou2022spes} --- which can prove query equivalence against \emph{all possible inputs}. 
They are most suitable for common queries (such as select-project-join) and give the highest possible guarantee; as a result, they are also more expensive. 
\slabcity uses full verification parsimoniously, only when our tester and bounded verifier are not able to disprove a query. 









\subsection{Performance Ranking}
\label{slab:sec:perfranking}

Our ultimate goal is to find equivalent queries that are \emph{faster}. 
We thus use a performance ranker to select a query with the best performance among the equivalent ones. 
In this paper, we use \texttt{EXPLAIN} cost estimates as a proxy for the actual performance. 
However, \slabcity can also use the actual query latency by running candidate queries against a small sample of the data (or even the entire data, when the overhead can be well-justified, e.g., for reporting dashboards where the same query will run many times). 
The choice of how to estimate query performance is an orthogonal question and is not essential to our synthesis-based algorithm. 
In Section \ref{slab:sec:ablation}, we present an ablation study to see the impact of using \texttt{EXPLAIN} cost estimates on the latency of the final synthesized queries. 
